\section{结论}
本文在SkillNet-demo上实现并检验了两项可运行成果：证据约束技能判别器，以及证据条件化的网络路径优化。
前者将可反驳的资源/接口证据与独立产物回执连接，输出修复、被测范围通过、证据不足和越界，
在构造留出诊断检出140/140个已知机械故障，并在288次自建程序输出上区分正确与缺陷实现。
自然包低覆盖同时表明，它尚不是通用技能质量认证模型。

后者在18节点公开数据派生任务网络中，把版本匹配、条件证据与完整兼容路径结合，
主留出三档平均成功70.83\%，高于原生元数据路径12.50\%；
对人工及全局风险路径仅提高4.17个百分点，宽预算成功数持平。
版本匹配降低了沿用旧代码证据的失败，完整路径减少了逐节点截断的能力缺口。
但前期取证使本批总调用成本高于人工流程，组合压力下只通过12/24。
因此成果是具有明确反例和适用条件的可复现方法及数据，而非对自然SkillNet网络普遍领先的主张。
